\section*{Chapitre \ref{ch:s2:zero-day-options-and-dealer-gamma-flows} --- Options à zéro jour et flux de gamma des dealers}

\begin{solution}{exo:s2:zero-day-options-and-dealer-gamma-flows:1}
Le delta de leurs options varie de $-5 \times 2 = -10$ unités d'indice~; pour rester couverts, ils achètent 10 unités, après la hausse.
\end{solution}

\begin{solution}{exo:s2:zero-day-options-and-dealer-gamma-flows:2}
Le gamma à la monnaie vaut $\varphi(d_1)/(S\sigma\sqrt\tau)$ avec $d_1 \to 0$~: il croît comme $1/\sqrt\tau$. Il reste peu de temps, un petit mouvement décide
si l'option finit dans la monnaie ou hors de la monnaie, si bien que son delta passe de près de 0 à près de 1.
\end{solution}

\begin{solution}{exo:s2:zero-day-options-and-dealer-gamma-flows:3}
Les strikes ronds sont espacés de 0.5 et la fenêtre couvre 0.05 de chaque côté, soit 0.1 sur chaque 0.5~: pour une clôture répartie uniformément sur un intervalle très supérieur à 0.5, la
part vaut $0.1/0.5 = 20~\%$.
\end{solution}

\begin{solution}{exo:s2:zero-day-options-and-dealer-gamma-flows:4}
La convention place les dealers en position longue sur tous les calls, situés au-dessus de l'ouverture, et en position courte sur tous les puts~: le gamma est négatif sous l'ouverture et
positif au-dessus par construction, de sorte qu'un changement de signe près de l'ouverture est inscrit dans l'hypothèse plutôt que trouvé dans les données.
\end{solution}

\begin{solution}{exo:s2:zero-day-options-and-dealer-gamma-flows:5}
Les côtés moyens sont les mêmes chaque jour~; ce qui détermine le gamma des dealers d'une journée, ce sont les côtés des clients ce jour-là, qui varient avec un écart-type de
0.5 autour de la moyenne. Les positions ouvertes n'apportent aucune information à leur sujet.
\end{solution}

\begin{solution}{exo:s2:zero-day-options-and-dealer-gamma-flows:6}
Les dealers en gamma court achètent après les hausses et vendent après les baisses, ce qui amplifie les mouvements (1.48~\% par jour contre 0.99~\% sans rétroaction)~; les dealers en gamma long
s'opposent aux mouvements (0.81~\%). Ni et ses coauteurs ont mis en évidence ce canal non informationnel sur les actions individuelles~: le rééquilibrage des couvertures a accru la volatilité et
la probabilité de grands mouvements.
\end{solution}

\begin{solution}{exo:s2:zero-day-options-and-dealer-gamma-flows:7}
Rien~: sans impact, l'indice est une marche aléatoire. La règle fondée sur le régime réel rapporte 0.03 point de base par jour ($t$ = 0.05). La transaction inconditionnelle affiche
$-1.38$ avec $t = -2.2$ sur ces 2\,000 jours, un rappel~: un $t$ de 2 apparaît par hasard dans une série dont la moyenne vraie est nulle.
\end{solution}

\begin{solution}{exo:s2:zero-day-options-and-dealer-gamma-flows:8}
Le signe estimé concorde avec la réalité sur seulement 56~\% des jours, et la transaction conditionnée par ce signe rapporte 1.02 point de base les jours de gamma court estimés
($t$ = 1.0), contre 3.12 les jours réels. L'énoncé traite une hypothèse sur les côtés comme une mesure.
\end{solution}

\begin{solution}{pb:s2:zero-day-options-and-dealer-gamma-flows:1}
\begin{enumerate}
\item Une option le jour de son échéance~; son gamma près du strike croît comme $1/\sqrt\tau$ à l'approche de la clôture.
\item 2.3 millions de contrats par jour, soit 59~\% du volume du SPX en 2025~; un record de 62.4~\% en août 2025.
\item Gamma court~: acheter après les hausses, vendre après les baisses~; gamma long~: l'inverse.
\item Le rééquilibrage des couvertures des teneurs de marché accroît la volatilité des actions et la probabilité de grands mouvements, sans information sous-jacente.
\item Gamma des dealers issu des positions ouvertes et d'un côté supposé~; le niveau où cette estimation change de signe.
\item Les dealers sont longs de tous les calls et courts de tous les puts~: les clients vendent des calls et achètent des puts.
\item Sur 56.0~\% des jours, avec une corrélation de 0.23.
\item La convention sur 96.0~\% des 200 premiers jours, la réalité sur 36.5~\%.
\item Un dixième du delta non couvert toutes les cinq minutes, le reste dans les 30 dernières minutes~; chaque transaction déplace la barre suivante.
\item 1.48~\% les jours de gamma court, 0.81~\% les jours de gamma long, 0.99~\% sans rétroaction.
\item Une clôture attirée vers un strike élevé par la couverture en gamma long~; non~: 21.4~\% des clôtures sont proches d'un strike rond les jours de gamma long, contre 20~\% par hasard.
\item Un gamma concentré sur un seul strike en fin de journée et une couverture qui suit de près le delta de chaque option.
\item \textbf{Résultat nommé.} L'estimation conventionnelle du gamma se trompe de signe sur 44~\% des jours (corrélation de 0.23 avec la réalité)~; la
transaction des 30 dernières minutes rapporte 3.12 points de base par jour les jours de gamma court réels et $-3.40$ les jours de gamma long, une règle utilisant le vrai côté atteint un ratio de Sharpe
de 1.68, et la même règle appliquée à l'estimation donne 0.40 avant coûts et $-0.11$ après coûts.
\item Le rendement des 30 dernières minutes est prédit par celui du reste de la journée, sur plus de 60 marchés à terme, un lien avec la couverture en gamma.
\item 0.27 point de base par jour, un ratio de Sharpe de 0.14.
\item La règle avec le vrai côté (1.16 après coûts)~; sur l'estimation, rien.
\item Des données de transactions avec indicateurs d'initiateur, des enquêtes auprès des dealers, ou des données de bourse sur les volumes des clients et des teneurs de marché.
\item Les positions ouvertes connues avant la séance, des coûts chaque jour, l'estimation du régime faite le matin même.
\item Gamma long sous le \omterm{def:s2:zero-day-options-and-dealer-gamma-flows:flip}{niveau de bascule du gamma} (gamma flip).
\item Sur la connaissance du côté des options que détiennent les dealers.
\end{enumerate}
\end{solution}

\begin{solution}{iq:s2:zero-day-options-and-dealer-gamma-flows:1}
Leurs options perdent du delta quand l'indice monte et en gagnent quand il baisse~; pour rester couverts, ils achètent après les hausses et
vendent après les baisses. Leurs transactions s'ajoutent au mouvement, d'autant plus que le gamma à zéro jour est élevé en fin de journée.
\end{solution}

\begin{solution}{iq:s2:zero-day-options-and-dealer-gamma-flows:2}
Sommer les positions ouvertes multipliées par le gamma, par strike, sous une hypothèse sur les côtés. La tester contre tout ce qui révèle la réalité~: volume signé par
type de participant, la volatilité réalisée et l'autocorrélation des rendements intrajournaliers par régime estimé, et la stabilité du signe de l'estimation.
\end{solution}

\begin{solution}{iq:s2:zero-day-options-and-dealer-gamma-flows:3}
Si les dealers sont longs de gamma à cet endroit, l'indice peut être retenu près de ce niveau~; s'ils sont courts, une sortie peut s'accélérer. Il faut trader de petites tailles, avec un risque défini (papillons
plutôt que straddles nus), et seulement si l'historique montre une \omterm{def:s2:zero-day-options-and-dealer-gamma-flows:pinning}{aimantation au strike} (pinning) au-delà du hasard.
\end{solution}

\begin{solution}{iq:s2:zero-day-options-and-dealer-gamma-flows:4}
Un gamma élevé près de la clôture~: un mouvement de quelques points modifie brutalement la couverture, et un gap ne peut pas être couvert. Il faut limiter par la concentration sur les strikes, par la perte
sur un mouvement de plusieurs points dans la dernière heure, et par la capacité de couverture près de la clôture.
\end{solution}

\begin{solution}{iq:s2:zero-day-options-and-dealer-gamma-flows:5}
Positions ouvertes issues du fichier du matin, volume intrajournalier par strike, cotations en direct pour les volatilités implicites, un gamma par strike sur une grille de niveaux d'indice, l'hypothèse sur les côtés
en configuration, et des alertes lorsque le signe de l'estimation ou le \omterm{def:s2:zero-day-options-and-dealer-gamma-flows:flip}{niveau de bascule du gamma} se déplace~; avec des contrôles sur les cotations périmées.
\end{solution}

\begin{solution}{iq:s2:zero-day-options-and-dealer-gamma-flows:6}
À la monnaie avec des taux nuls, $d_1 = \tfrac12\sigma\sqrt\tau$, et $\Gamma = \varphi(d_1)/(S\sigma\sqrt\tau)$. Quand $\tau \to 0$, $\varphi(d_1) \to
1/\sqrt{2\pi}$, donc $\Gamma \approx 1/(S\sigma\sqrt{2\pi\tau})$~: il croît comme $\tau^{-1/2}$.
\end{solution}
